\section{Fazit und Ausblick}
\label{sec:conclusion}

\subsection{Inhaltliche Stichpunkte}
\begin{itemize}
  \item \textbf{Kernfrage beantworten}
  \begin{itemize}
    \item Im Notebook~07 erzielt Model~2 mit Full Context im betrachteten
          240\,000-Zeilen-Testlauf höhere Gesamtwerte als die Baseline:
          Macro-F1 0{,}953685 statt 0{,}911805, ROC-AUC 0{,}994932 statt
          0{,}986353.
    \item Vorsichtige Schlussfolgerung: In dieser Evaluation verbessert
          sich die beobachtete Erkennungsleistung. Eine statistisch
          signifikante Verbesserung ist ohne Signifikanztest nicht belegt.
  \end{itemize}
  \item \textbf{Bedeutung der Methodik}
  \begin{itemize}
    \item Die hybride Pipeline kombiniert reale Messreihen mit
          klassenspezifischen TabGAN- und DAE-Anteilen.
    \item Im separaten Generatorvergleich erzielt der Combined-Nodes-Hybrid
          Macro-F1 0{,}8890 und ROC-AUC 0{,}9845 auf dem Original-Testdatensatz.
    \item Empirische Begründung der Generatorwahl und unabhängige
          Testdaten als zentrale Voraussetzungen der Aussagekraft hervorheben.
  \end{itemize}
  \item \textbf{Grenzen und zukünftige Arbeit}
  \begin{itemize}
    \item Der aktuelle Aufbau umfasst acht simulierte Knoten; Skalierung auf
          größere Netzwerke untersuchen.
    \item Physische Ladesäulen und echte Kommunikationsbedingungen sind
          noch zu evaluieren.
    \item Latenz, Bandbreite, Ausfälle, Datenschutz und manipulierte
          Peer-Berichte als offene Fragestellungen aufnehmen.
    \item Fehlende Signifikanzanalyse und noch zu berechnende
          Falsch-Positiv-Rate als Grenzen der vorliegenden Auswertung benennen.
  \end{itemize}
\end{itemize}

\subsection{Noch zu erklären}
\begin{itemize}
  \item Die Forschungsfrage anhand des vollständigen, einheitlichen
        Hauptvergleichs beantworten und dabei keine statistische
        Signifikanz behaupten, solange sie nicht getestet wurde.
  \item Abgrenzung der zwei Evidenzstränge: Generatorvergleich auf einem
        separaten Original-Testdatensatz und P2P-Modellvergleich auf den
        zusammengeführten Knotentestdaten.
  \item Grenzen der Übertragbarkeit von synthetischen/simulierten Daten auf
        reale EVSE und Einsatzbedingungen erläutern.
\end{itemize}

\subsection{Grafikbedarf}
\begin{itemize}
  \item Keine neue Ergebniskurve nötig; auf die Hauptvergleichstabelle
        und die kompakte Metrikgrafik in Kapitel~\ref{sec:results} verweisen.
  \item Optional eine abschließende Übersicht mit Forschungsfrage,
        Kernergebnis, Einschränkungen und nächsten Schritten.
\end{itemize}
